\section{Marco Teórico}

El análisis del impacto del salario mínimo 
en mercados laborales con alta informalidad, 
como el de Yucatán, se sustenta principalmente 
en dos enfoques de la teoría económica moderna: 
el modelo de competencia perfecta y el modelo 
de búsqueda (search models).

\subsection{El Modelo de Sector Dual y el Efecto Desplazamiento}
De acuerdo con la teoría neoclásica tradicional, 
en un mercado laboral segmentado, el sector formal
está sujeto a regulaciones institucionales (salario mínimo), 
mientras que el sector informal opera bajo condiciones de equilibrio de mercado. 

Al incrementar el salario mínimo por encima de la 
productividad marginal del trabajo en el sector formal, 
se genera un exceso de oferta laboral. Las empresas 
formales, al no poder absorber este costo, reducen 
su demanda de mano de obra, provocando un 
\textit{efecto de desplazamiento} hacia el sector 
informal, donde los salarios son flexibles y no 
existe una barrera de entrada legal.

\subsection{Teoría de las Barreras de Entrada y Rigidez Laboral}
Complementariamente, la teoría de los modelos de 
búsqueda (\textit{Search and Matching Models}) sugiere 
que el salario mínimo actúa como un \textbf{precio piso} 
que eleva el estándar de contratación. 

Para los nuevos entrantes al mercado laboral de 
Yucatán (jóvenes y trabajadores con baja escolaridad), 
un salario mínimo elevado puede representar una barrera 
de entrada. Si el costo de contratación supera la 
productividad esperada del trabajador, la firma 
optará por no abrir la vacante formal o recurrir a 
esquemas de contratación informal para evitar los 
costos de seguridad social, validando la hipótesis 
de \textit{ajuste de flujo} planteada en este documento.

\subsection{Efecto Faro (Signaling Effect)}
Es relevante mencionar el ''Efecto Faro'', donde los 
incrementos en el salario mínimo formal sirven 
como referencia para la fijación de precios y 
salarios en el sector informal. En Yucatán, este 
fenómeno puede generar presiones inflacionarias 
locales o reajustes en la estructura de costos de 
los micronegocios, influyendo indirectamente en 
la persistencia de la informalidad.

\subsection{Productividad Marginal y Heterogeneidad Regional}
La literatura económica reciente subraya que el 
impacto del salario mínimo no es uniforme y depende 
críticamente de la \textbf{productividad marginal del trabajo} 
preexistente en cada región \citep{Neumark2021}. 
En estados con sectores industriales de alta 
capitalización, como el norte de México, la 
productividad suele situarse holgadamente por encima 
del salario mínimo, lo que permite absorber los 
incrementos sin alterar significativamente la demanda 
de empleo formal.
Por el contrario, en regiones como el sureste 
mexicano, donde predomina el sector servicios y 
el comercio minorista de baja escala, la brecha 
entre la productividad laboral y el costo legal del 
trabajo es estrecha. De acuerdo con estudios del 
\textbf{Banco de México} \citep{Banxico2022}, 
incrementos acelerados en el salario mínimo en zonas 
de baja productividad pueden incentivar el tránsito 
hacia la informalidad como un mecanismo de ajuste 
de costos para las micro y pequeñas empresas (MiPyMEs). 

Bajo esta óptica, el salario mínimo 
actúa como un ``impuesto a la formalidad'' 
en Yucatán, donde la estructura productiva 
local —caracterizada por una menor densidad de 
capital por trabajador— dificulta la absorción 
de costos laborales crecientes, validando la 
preocupación de que la política salarial nacional 
puede profundizar las brechas de desarrollo regional.